\documentclass[../../../main.tex]{subfiles}
\begin{document}
\subsection{Ergotropy of Qubit}

\subsubsection{General formula.}
The ergotropy of a qubit in state $\rho$ is given by
\begin{equation*}
  \mathcal{E} [\rho]
  = B(2\rho_{11}- 1 + \sqrt{1-4 \det \rho})
\end{equation*}
Assuming that the higher level energy have higher population, that is $\rho_{11}>0.5$, we have ergotropy due to population inversion
\begin{equation*}
  \mathcal{E}_\text{pop}(\rho)
  = 2B(\rho_{11}-\rho_{00}) = 2B(2\rho_{11}-1)
\end{equation*}
Then the ergotropy associated with coherence is
\begin{equation*}
  \mathcal{E}_\text{coh}[\rho]
  = B \left[  1 - 2 \rho_{11} + \sqrt{(1-2\rho_{11})^2+ 4|\rho_{01}|^2}  \right]
\end{equation*}
This value is basically how much ergotropy is lost when I remove the coherence while keeping the populations unchanged.
Making the state diagonal, in other words.
Note the explicit dependence on diagonal element of this ergotropy, the coherence ergotropy is zero if $\rho_{01}=0$.
Writing the state as
\begin{equation*}
  \rho
  =
  \frac12\left(\mathbf{1}+\mathbf r\cdot\boldsymbol{\sigma}\right)
  =
  \frac12
  \begin{bmatrix}
    1+z  & x-iy \\
    x+iy & 1-z
  \end{bmatrix}
\end{equation*}
makes the geometric interpretation of coherence very clear: coherence is associated with $xy$ place projection.

\subsubsection{Pure and mixed state.}
Suppose we are given the pure and mixed state of spin chain in the form
\begin{gather*}
  \ket{\psi(0)} = \left(\alpha\ket{0}_1 + \beta\ket{1}_1\right)\bigotimes_{i=2}^{N}\ket{0}_i  = \alpha\ket{\mathbf{0}} + \beta\ket{\mathbf{1}}\\
  \rho(0) = \left(q\ket{1}\bra{1}_1 + (1-q)\ket{0}\bra{0}_1\right)\bigotimes_{i=2}^{N}\ket{0}\bra{0}_i = q\ket{\mathbf{1}}\bra{\mathbf{1}} + (1-q)\ket{\mathbf{0}}\bra{\mathbf{0}}
\end{gather*}
with $q = (1+|\beta^2|)/2$.
The first lines written in such way to demonstrate that the first site of spin chain is initially pure and mixed respectively.
Also, we do not simply set $q=|\beta|^2$ so that, with the same choice of $\alpha$ and $\beta$, the initial state of both pure state and mixed state encode the same initial ergotropy.
Alternatively, the state can be written as
\begin{equation*}
  \rho_\text{coh}
  = \begin{bmatrix}
    \alpha^2    & \alpha\beta^* \\
    \alpha\beta & |\beta|^2
  \end{bmatrix}
  \qquad
  \rho_\text{mix}
  = \begin{bmatrix}
    1-q & 0 \\
    0   & q
  \end{bmatrix}
\end{equation*}
As usual the pure state exist on the surface of Bloch sphere, while the mixed state on the interior, specifically on $z$ axis with $z = 1-2q$.
First note that the local Hamiltonian in this case is
\begin{equation*}
  H = \begin{bmatrix}
    -B & 0 \\
    0  & B
  \end{bmatrix}
\end{equation*}
The initial ergotropy for both cases are
\begin{equation*}
  \mathcal{E}\left[ \rho_{\text{pure}}^{(1)} \right]  = 2B|\beta|^2 \qquad
  \mathcal{E}\left[ \rho_{\text{mix}}^{(1)} \right] = 2B (2q-1)
\end{equation*}
As for the transferred ergotropy
\begin{equation*}
  \mathcal{E}\left[ \rho_{\text{pure}}^{(N)} \right]= 2|f_N(t)|^2\sin^2\left(\frac{\theta}{2}\right)-1+\sqrt{1+4\sin^4\left(\frac{\theta}{2}\right)|f_N(t)|^2\left(|f_N(t)|^2-1\right)}
\end{equation*}
and
\begin{equation*}
  \mathcal{E}\left[ \rho_{\text{mix}}^{(N)} \right] = 2[2q |f_N|^2 - 1]
\end{equation*}
Both still dependent on $\theta$, so by averaging it over the Bloch sphere we can have a quantity that characterize the quantum channel as a channel for ergotropy transport
\begin{equation*}
  \braket{\mathcal{E}_\text{coh}} = |f_N|^2-1 + \frac{\arcsin 2|f_N|\sqrt{1-|f_N|^2 }}{4|f_N|\sqrt{1-|f_N|^2}} + \frac{1 }{2} |1-2|f_N|^2 |
\end{equation*}

\subsubsection{Arbitrary qubit state.}
Now suppose we create an ensemble the pure state
\begin{equation*}
  \ket{\psi_i}
  =
  \cos\frac{\theta_i}{2}\ket{0}
  +
  e^{i\phi_i}\sin\frac{\theta_i}{2}\ket{1},
\end{equation*}
The expression for this state is
\begin{equation*}
  \rho^{(1)}  =
  \begin{bmatrix}
    p\cos^2\frac{\theta_i}{2}
    +(1-p)\cos^2\frac{\theta_j}{2}
     &
    p e^{-i\phi_i}\frac{\sin \theta_i}{2}
    +(1-p)e^{-i\phi_j}\frac{\sin \theta_j}{2}
    \\\\
    p e^{i\phi_i}\frac{\sin \theta_i}{2}
    +(1-p)e^{i\phi_j}\frac{\sin \theta_j}{2}
     &
    p\sin^2\frac{\theta_i}{2}
    +(1-p)\sin^2\frac{\theta_j}{2}
  \end{bmatrix}
\end{equation*}
To simplify our expression, we can use the following parametrization: \((\theta_j,\phi_j)=(-\theta_i,-\phi_i)\), \((\theta_i,-\phi_i)\), and \((-\theta_i,\phi_i)\); respectively we define these constraints as $C_2$, $C_3$, and $C_4$.
The constraints are sufficient to parameterize the complete single-qubit state space.
By varying the parameters $(\theta,\phi,p)$, every physical qubit state within and on the surface of the Bloch sphere can be represented under each constraint.
Note that negative \(\theta\) should be interpreted as a parametrization device, rather than as a conventional polar angle, since $0 \leq \theta \leq \pi$.
Equivalently, the negative $\theta$ constraint can be represented by $(\theta_2=\theta_1,\phi_2 = \pi-\phi_1)$ and $(\theta_2=\theta_1,\phi_2 = \pi+\phi_1)$.

\begin{figure}
  \centering
  \begin{tikzpicture}[tdplot_main_coords,scale=4]
    \draw[thick,->] (0,0,0) -- (1.2,0,0) node[anchor=north east]{$x$};
    \draw[thick,->] (0,0,0) -- (0,1.1,0) node[anchor=north west]{$y$};

    \def\r{1}
    \def\th{60}
    \def\thm{-\th}
    \def\ph{35}
    \def\phm{-\ph}
    \pgfmathsetmacro{\rxy }{sin(\th)}
    \pgfmathsetmacro{\rxym}{sin(\thm)}

    \coordinate (O) at (0,0,0);

    \tdplotsetcoord{A}{\r}{\th}{\ph}
    \tdplotsetcoord{A1}{\r}{\thm}{\phm}
    \tdplotsetcoord{A2}{\r}{\th}{\phm}
    \tdplotsetcoord{A3}{\r}{\thm}{\ph}
    \tdplotsetcoord{OO}{\r}{0}{0}

    \draw[thick,->,>=Stealth] (O) -- (A) node [anchor=south east]{$\ket{\psi_1}$};
    \draw[thick,dashed,->,>=Stealth] (O) -- (A1) node [anchor=south west]{$\ket{\psi_2}$};
    \draw[thick,->,>=Stealth] (O) -- (A2) node [anchor=south west]{$\ket{\psi_3}$};
    \draw[thick,dashed,->,>=Stealth] (O) -- (A3) node [anchor=north east]{$\ket{\psi_4}$};
    \draw[thick,->,>=Stealth] (O) -- (OO) node [anchor=south]{$\ket{0}$};

    \tdplotdrawarc[->]{(O)}{\rxy}{0}{\ph}{anchor=south}{$\phi$}
    \tdplotsetthetaplanecoords{\ph}
    \tdplotdrawarc[->,tdplot_rotated_coords]{(0,0,0)}{\r}{0}{\th}{anchor=north west}{$\theta$}

    \draw[dashed]  (O) -- (Axy);
    \draw[dashed]  (A) -- (Axy);
    \draw[dashed]  (O) -- (A1xy);
    \draw[dashed]  (A1) -- (A1xy);
    \draw[dashed]  (O) -- (A2xy);
    \draw[dashed]  (A2) -- (A2xy);
    \draw[dashed]  (O) -- (A3xy);
    \draw[dashed]  (A3) -- (A3xy);

    \def\alp{20}
    \tdplotsetrotatedcoords{\alp}{0}{0}
    \begin{scope}[tdplot_rotated_coords]
      \draw[thick] (0,-\r) arc (-90:90:\r);
      \draw[dashed] (0,\r) arc (90:270:\r);
    \end{scope}

    \tdplotsetrotatedcoords{\alp}{90}{0}
    \begin{scope}[tdplot_rotated_coords]
      \draw[thick] (0,\r) arc (90:270:\r);
      % \draw[thick] (0,\r) arc (90:360+90:\r);
    \end{scope}

  \end{tikzpicture}
  \caption*{
    Figure: Bloch sphere representation of mixed states created from pure state ensemble.
    The states \(\ket{\psi_2}\), \(\ket{\psi_3}\), and \(\ket{\psi_4}\) correspond to the coordinate constraints \((\theta,\phi)=(-\theta_1,-\phi_1)\), \((\theta_1,-\phi_1)\), and \((-\theta_1,\phi_1)\), respectively.
  }
\end{figure}

To avoid misinterpretation, we therefore choose the third constrain \((\theta_1,-\phi_1)\).
The state in this case is given by
\begin{equation*}
  \rho ^{(1)} =
  \begin{bmatrix}
    \cos ^2 \theta/2                                                 & \frac{\sin \theta}{2} \left[-i2p \sin \phi + e^{i\phi}  \right] \\
    \frac{\sin \theta  }{2} \left[i2p \sin\phi + e^{-i\phi}  \right] & \sin ^2 \frac{\theta  }{2}
  \end{bmatrix}
\end{equation*}
The last site then will carry the information from the first state and evolve as
\begin{equation*}
  \rho ^{(N)}  =
  \begin{bmatrix}
    1- \sin^2 \theta/2 |f_N|^2                                             & \frac{\sin \theta}{2} \left[-i2p \sin \phi + e^{i\phi}  \right]f_N \\
    \frac{\sin \theta  }{2} \left[i2p \sin\phi + e^{-i\phi}  \right]f_N ^* & \sin ^2 \frac{\theta  }{2}|f_N|^2
  \end{bmatrix}
\end{equation*}

\subsubsection{Specific state.}
We can choose specific state such that the coherence and population have the same contribution.
This family of state is given by
\begin{gather*}
  p = \frac{1 }{2 }- \frac{1 }{2} \sqrt{1 - \frac{1-9\cos^2\theta}{\sin^2\theta\sin^2\phi}}\\
  \pi/2<\theta\le \pi - \theta_0 \qquad
  \phi = \phi_0 + \lambda( \pi - 2\phi_0)
\end{gather*}
with $\theta_0 = \arccos(1/3)$, $\phi_0 = \arccos(\sqrt{ 1-8 \cot ^2 \theta}) $, and $0 \le \lambda \le 1$.
For convenience’s sake, take $\lambda=1/2$ to obtain $\phi=\pi/2$ and simplify the expressions into
\begin{equation*}
  p=\frac12+\sqrt2\cot\theta\qquad
  \frac{\pi}{2}<\theta\le\pi-\theta_0.
\end{equation*}

These conditions are generated by the equal contribution requirement
\begin{equation*}
  p^2-p=- \frac{1-9\cos^2\theta}{4\sin^2\theta\sin^2\phi}
\end{equation*}
Notice that there are three independent parameters $(p,\theta,\phi)$ and one single constrain.
Hence, the system has $N_\text{DOF} = 3-1 =2$ degree of freedom and a single parameter cannot determine all three variables.
Imposing additional constrain $\phi = \pi/2$, we can express both $p$ and $\phi$ in terms of $\theta$.
Also do note that the population inversion requirement $\theta>\pi/2$ does not count as constrain since it only restrict the domain of $\theta$.

Ignoring extra parameter $\lambda$, our state is given by
\begin{gather*}
  p = \frac{1 }{2 }- \frac{1 }{2} \sqrt{1 - \frac{1-9\cos^2\theta}{\sin^2\theta\sin^2\phi}}\\
  \pi/2<\theta\le \pi - \theta_0 \qquad
  \phi \in[\phi_0,\pi-\phi_0] \cup[\pi+\phi_0,2\pi-\phi_0]
\end{gather*}
Notice that the second branch is naturally paired under $(\theta_1,-\phi_1)$ construction.

\subsubsection{Arbitrary qubit.}
The full ergotropy is
\begin{align*}
  \mathcal{E} \left[ \rho^{(1)} \right]
   & = B \left( 2\sin^2 \frac{\theta }{2} - 1 + \sqrt{1 - 4p(1-p)\sin ^2\theta \sin^2 \phi}  \right)
\end{align*}
and the final site
\begin{multline*}
  \mathcal{E} \left[ \rho^{(N)} \right]     = B \bigg( 2\sin^2 \frac{\theta }{2} |f_N|^2 - 1
  \\  \qquad+  \sqrt{1 - 4|f_N|^2 \left[ p(1-p)\sin ^2\theta \sin^2 \phi + \sin ^4 \frac{\theta }{2 }(1-|f_N|^2) \right]}  \bigg)
\end{multline*}
The population ergotropy for both
\begin{equation*}
  \mathcal{E}_\text{pop} \left[ \rho^{(1)} \right]= 2B\max(2 \sin^2 \frac{\theta }{2 }-1,0)  \quad
  \mathcal{E}_\text{pop} \left[ \rho^{(N)} \right]= 2B\max(2 |f_N|^2\sin ^2\frac{\theta}{2 }-1,0)
\end{equation*}
And the coherence ergotropy
\begin{align*}
  \mathcal{E}_\text{coh}\left[ \rho^{(1)}\right]
   & = B \left[  1 - 2 \sin ^2 \frac{\theta }{2} + \sqrt{1-4p(1-p)\sin ^2\theta \sin^2 \phi}  \right]
\end{align*}
and
\begin{multline*}
  \mathcal{E}_\text{coh}\left[ \rho^{(N)}\right]
  = B \bigg[  1 - 2 |f_N|^2\sin ^2 \frac{\theta }{2} \\
    + \sqrt{1-4 |f_N|^2 \left[ p(1-p)\sin ^2\theta \sin^2 \phi
        + \sin ^4 \frac{\theta }{2 }(1-|f_N|^2) \right] }  \bigg]
\end{multline*}

\subsubsection{Derivation: general ergotropy and purity.}
Recall the state evolution $\rho^{(1)}(0) \rightarrow \rho^{(N)}(t)$
\begin{equation*}
  \begin{bmatrix}
    1 -\rho_{11}^{(1)}  & \rho_{01}^{(1)} \\
    [\rho_{01}^{(1)}]^* & \rho_{11}^{(1)}
  \end{bmatrix}
  \rightarrow
  \begin{bmatrix}
    1 -\rho_{11}^{(1)}|f_N(t)|^2 & \rho_{01}^{(1)}f_N(t)     \\
    [\rho_{01}^{(1)}f_N(t)]^*    & \rho_{11}^{(1)}|f_N(t)|^2
  \end{bmatrix}
\end{equation*}
The eigenvalue of this initial state is
\begin{align*}
  \det \begin{bmatrix}\rho_{00}-\lambda & \rho_{01}\\ \rho_{10} & \rho_{11}-\lambda\end{bmatrix} & = \lambda^2-(\rho_{00}+\rho_{11})\lambda+(\rho_{00}\rho_{11}-\rho_{01}\rho_{10}) \\
                                                                                                 & = \lambda^2 - \mathrm{Tr} [\rho]\lambda+\det \rho                                \\
  0                                                                                              & =  \lambda^2 -\lambda + \det \rho
\end{align*}
and using quadratic formula
\begin{equation*}
  \lambda_{\pm} = \frac{1 \pm \sqrt{1-4 \det \rho}}{2 }
\end{equation*}
From this the passive energy is
\begin{equation*}
  E_\mathrm{st} [\rho]= \lambda_+ H_{00} + \lambda_- H_{11}  = (\lambda_--\lambda_+)B = - B \sqrt{1-4 \det \rho}
\end{equation*}
The energy stored is simply
\begin{equation*}
  E_\mathrm{pas} [\rho]= \mathrm{Tr}  \begin{bmatrix}
    \rho_{00} & \rho_{01} \\
    \rho_{10} & \rho_{11}
  \end{bmatrix}
  \begin{bmatrix}
    -B & 0 \\
    0  & B
  \end{bmatrix}
  = \mathrm{Tr}
  \begin{bmatrix}
    -B\rho_{00}  & B\rho_{01} \\
    -B \rho_{10} & B\rho_{11}
  \end{bmatrix}
  = B(\rho_{11}-\rho_{00})
\end{equation*}
Therefore the ergotropy
\begin{equation*}
  \mathcal{E} [\rho]
  = B(\rho_{11}-\rho_{00} + \sqrt{1-4 \det \rho})
  = B(2\rho_{11}- 1 + \sqrt{1-4 \det \rho})
\end{equation*}

Also consider the determinant,
\begin{align*}
  \det\rho     & =  \frac14\left[(1+z)(1-z)-(x-iy)(x+iy)\right]        \\
  \det\rho     & =  \frac14(1-x^2-y^2-z^2) = \frac14(1-|\mathbf{r}|^2) \\
  |\mathbf{r}| & =  \sqrt{1-4\det\rho}
\end{align*}
we can write
\begin{equation*}
  \mathcal{E}[\rho] = B(|\mathbf{r}|-z)
\end{equation*}
We can also obtain the same result by considering the stored and passive energy in terms of Bloch vector first.
For the stored energy
\begin{equation*}
  E_{\rm st} = \mathrm{Tr}(\rho(-B\sigma_z)) = -B\mathrm{Tr}(\rho\sigma_z)  = -Bz
\end{equation*}
As for the passive, we express the eigenvalue as $\lambda = (1\pm |\mathbf{r}|)/2$ and construct the passive state as
\begin{equation*}
  \rho_{\rm pas} = \begin{pmatrix} \lambda_+&0\\ 0&\lambda_- \end{pmatrix} = \frac12 \begin{pmatrix} 1+r&0\\ 0&1-r \end{pmatrix}
\end{equation*}
such that the passive energy is
\begin{equation*}
  E_{\rm pas}
  = (-B)\lambda_+ + B\lambda_-
  \frac{-B-Br+B-Br}{2} =-Br
\end{equation*}
and the ergotropy
\begin{equation*}
  \mathcal{E} = B(r-z)
\end{equation*}
With ground state energy as $E_0 = -B$, we can write also both as $ E_{\rm st}=E_0z$ and  $E_{\rm pas}=E_0r$, such that $\mathcal{E} = E_0(z-r)$.


For the case of population inversion, by assuming that $\rho_{11}\geq 0.5$, the derivation is very simple
\begin{align*}
  \mathcal{E}_\text{pop}(\rho)
   & = B(\rho_{11}-\rho_{00}) -  B(\rho_{00}-\rho_{11}) \\
   & = B(\rho_{11}-\rho_{00}) +  B(\rho_{11}-\rho_{00}) \\
   & = 2B(\rho_{11}-\rho_{00}) = 2B(2\rho_{11}-1)
\end{align*}
Alternatively,  we can use our previous equation
\begin{align*}
  \mathcal{E}_\text{pop} [\rho]
   & = B \left[ 2\rho_{11}-1+ \sqrt{1-4(1-\rho_{11})\rho_{11}} \right] \\
   & = B \left[ 2\rho_{11}-1+ \sqrt{(1 - 2\rho_{11})^2 } \right]       \\
  \mathcal{E}_\text{pop} [\rho]
   & = B \left[ 2\rho_{11}-1+ |1 - 2\rho_{11}|\right]
\end{align*}
And by using $x+|x| = 2 \max (x,0)$ identify, we can obtain the most general formula without assuming $\rho_{11}>0.5$
\begin{equation*}
  \mathcal{E}_\text{pop} [\rho]  = 2B \max \left[ 2\rho_{11}-1,0 \right]
\end{equation*}

The coherence-associated ergotropy from coherence can be calculated
\begin{align*}
  \mathcal{E}_\text{coh}[\rho]
   & =  \mathcal{E}[\rho] -   \mathcal{E}_\text{pop}[\rho]                                      \\
   & = B \left[ 2 \rho_{11} - 1 + \sqrt{1-4 \det \rho} - 4\rho_{11} + 2  \right]                \\
   & = B \left[  1 - 2 \rho_{11} + \sqrt{1-4 (\rho_{11}(1-\rho_{11}) - |\rho_{01}|^2)}  \right] \\
  \mathcal{E}_\text{coh}[\rho]
   & = B \left[  1 - 2 \rho_{11} + \sqrt{(1-2\rho_{11})^2+ 4|\rho_{01}|^2}  \right]
\end{align*}

Next consider the determinant of initial and final state.
The initial state determinant
\begin{align*}
  \det\rho^{(1)}
   & = \rho_{11}^{(1)}(1-\rho_{11}^{(1)}) - |\rho_{01}^{(1)}|^2
  = \rho_{11}^{(1)} -\rho_{11}^{(1)}{}^2- |\rho_{01}^{(1)}|^2
\end{align*}
can be used to determine the evolved state
\begin{align*}
  \det \rho^{(N)} & = (1-\rho^{(1)}_{11}|f_N|^2)\rho^{(1)}_{11}|f_N|^2 - |\rho^{(1)}_{01}f_N|^2                                               \\
                  & = |f_N|^2 (\rho^{(1)}_{11}-{\rho^{(1)}_{11}}{}^2|f_N|^2-|\rho^{(1)}_{01}|^2)                                              \\
                  & = |f_N|^2(\rho^{(1)}_{11} - \rho^{(1)}_{11}{}^2-|\rho^{(1)}_{01}|^2 + \rho^{(1)}_{11}{}^2 - {\rho^{(1)}_{11}}{}^2|f_N|^2) \\
  \det \rho^{(N)} &
  = |f_N|^2[\det \rho^{(1)}+ \rho^{(1)}_{11}{}^2(1-|f_N|^2)]
\end{align*}

The general expression for $2 \times 2$ density matrix purity is
\begin{align*}
  \mathcal{P}(\rho)
   & =
  \mathrm{Tr}
  \begin{bmatrix}
    \rho_{00} & \rho_{01} \\
    \rho_{10} & \rho_{11}
  \end{bmatrix}
  \begin{bmatrix}
    \rho_{00} & \rho_{01} \\
    \rho_{10} & \rho_{11}
  \end{bmatrix}
  \\
   & =
  \mathrm{Tr}
  \begin{bmatrix}
    \rho_{00}^2+\rho_{01}\rho_{10}
     &
    \rho_{00}\rho_{01}+\rho_{01}\rho_{11}
    \\
    \rho_{10}\rho_{00}+\rho_{11}\rho_{10}
     &
    \rho_{10}\rho_{01}+\rho_{11}^2
  \end{bmatrix}
  \\
  \mathcal{P}(\rho)
   & =
  \rho_{00}^2+\rho_{11}^2
  +2|\rho_{01}|^2
\end{align*}
Using $\lambda_\pm = (1\pm \sqrt{1- \det \rho})/2$ we have
\begin{equation*}
  \lambda_+ + \lambda_-  = \mathrm{Tr} \rho =1\qquad
  \lambda_+\lambda_- = \det \rho
\end{equation*}
Hence we can express the trace in terms of determinant
\begin{equation*}
  \mathcal{P}(\rho) = (\lambda_1+\lambda_2)^2-2\lambda_1\lambda_2 = 1-2\det\rho
\end{equation*}
Alternatively, since $\det \rho = \rho_{00}\rho_{11}-|\rho_{01}|^2$ we can also obtain the same result
\begin{align*}
  \mathcal{P}(\rho) & =  (\rho_{00}+\rho_{11})^2  - 2\rho_{00}\rho_{11} + 2|\rho_{01}|^2 \\
  \mathcal{P}(\rho) & =1-2\left(\rho_{00}\rho_{11}-|\rho_{01}|^2\right)=1-2\det\rho
\end{align*}
The purity of the first state is trivial so for the final state
\begin{equation*}
  \mathcal{P} [\rho^{(N)}] =
  1-  2 |f_N|^2[\det \rho^{(1)}+ \rho^{(1)}_{11}{}^2(1-|f_N|^2)]
\end{equation*}
Then the purity difference can be expressed as follows
\begin{align*}
  \Delta\mathcal{P}[\rho^{(N)},\rho^{(1)}]
   & =
  -2|f_N|^2\det\rho^{(1)}
  -2|f_N|^2(\rho_{11}^{(1)})^2(1-|f_N|^2)
  +2\det\rho^{(1)}
  \\
   & =
  2\det\rho^{(1)}(1-|f_N|^2)
  -2|f_N|^2(\rho_{11}^{(1)})^2(1-|f_N|^2)
  \\
   & =
  2(1-|f_N|^2)
  \left[
    \det\rho^{(1)}
    -|f_N|^2(\rho_{11}^{(1)})^2
    \right]
\end{align*}

\subsubsection{Derivation: ergotropy for pure and mixed state.}
The initial state can be written respectively as
\begin{equation*}
  \rho_\text{coh}
  = \begin{bmatrix}
    \alpha^2    & \alpha\beta^* \\
    \alpha\beta & |\beta|^2
  \end{bmatrix}
  \qquad
  \rho_\text{mix}
  = \begin{bmatrix}
    1-q & 0 \\
    0   & q
  \end{bmatrix}
\end{equation*}
respectively.
The initial encoded ergotropy for pure state is
\begin{equation*}
  \mathcal{E}[\rho_\text{coh}]
  = B(2|\beta|^2-1 \sqrt{1-4(\alpha^2|\beta|^2-\alpha^2\beta^2)})
  = 2B|\beta|^2
\end{equation*}
Whereas the mixed state
\begin{equation*}
  \mathcal{E}[\rho_\text{mix}] = 2B(2q- 1)
\end{equation*}
Suppose both case has the same initial ergotropy, then $q = (1+|\beta|^2)/2$.

Next consider the transferred ergotropy.
Let us then write the density matrix in the following form
\begin{equation*}
  \rho_\text{coh}^{(N)}
  = \begin{bmatrix}
    \alpha^2 + |\beta|^2 (1-|f_N|^2) & \alpha\beta^*  f_N^* \\
    \alpha\beta  f_N                 & |\beta|^2 |f_N|^2
  \end{bmatrix}
\end{equation*}
and for the mixed state
\begin{equation*}
  \rho_\text{mix}^{(N)}
  = \begin{bmatrix}
    q + (1-q) (1-|f_N|^2) & 0            \\
    0                     & (1-q)|f_N|^2
  \end{bmatrix}
\end{equation*}
Next finding the determinant.
For the coherent
\begin{align*}
  \det \rho_\text{coh} & = [\alpha^2 + |\beta|^2 (1-|f_N|^2)]|\beta|^2|f_N|^2 - |\alpha|^2|\beta|^2|f_N|^2 \\
  \det \rho_\text{coh} & = |\beta|^4|f_N|^2 (1-|f_N|^2)
\end{align*}
and for mixed state we'll just use population inversion like usual.
Hence, the ergotropy for pure state
\begin{equation*}
  \mathcal{E}_\text{coh} (t) = 2|f_N(t)|^2\sin^2\left(\frac{\theta}{2}\right)-1+\sqrt{1+4\sin^4\left(\frac{\theta}{2}\right)|f_N(t)|^2\left(|f_N(t)|^2-1\right)}
\end{equation*}
and for mixed
\begin{equation*}
  \mathcal{E}_\text{mix}(t) = 2B((1-q)|f_N|^2-1)
\end{equation*}


\subsubsection{Derivation: general mixed state expression and constraint.}
Now consider a mixed state created from an ensemble of pure state
The pure state is given by
\begin{equation*}
  \ket{\psi_i}
  =
  \cos\frac{\theta_i}{2}\ket{0}
  +
  e^{i\phi_i}\sin\frac{\theta_i}{2}\ket{1},
\end{equation*}
The statistical mixed state with probabilities \(p\) and \(1-p\) is
\begin{align*}
  \rho^{(1)}
             & =
  p\ket{\psi_i}\bra{\psi_i}
  +
  (1-p)\ket{\psi_j}\bra{\psi_j}                              \\
             & =
  p
  \begin{bmatrix}
    \cos^2\frac{\theta_i}{2}
     &
    e^{-i\phi_i}\cos\frac{\theta_i}{2}\sin\frac{\theta_i}{2} \\\\
    e^{i\phi_i}\cos\frac{\theta_i}{2}\sin\frac{\theta_i}{2}
     &
    \sin^2\frac{\theta_i}{2}
  \end{bmatrix} \\
             & \qquad
  +(1-p)
  \begin{bmatrix}
    \cos^2\frac{\theta_j}{2}
     &
    e^{-i\phi_j}\cos\frac{\theta_j}{2}\sin\frac{\theta_j}{2} \\\\
    e^{i\phi_j}\cos\frac{\theta_j}{2}\sin\frac{\theta_j}{2}
     &
    \sin^2\frac{\theta_j}{2}
  \end{bmatrix} \\
  \rho^{(1)} & =
  \begin{bmatrix}
    p\cos^2\frac{\theta_i}{2}
    +(1-p)\cos^2\frac{\theta_j}{2}
     &
    p e^{-i\phi_i}\frac{\sin \theta_i}{2}
    +(1-p)e^{-i\phi_j}\frac{\sin \theta_j}{2}
    \\\\
    p e^{i\phi_i}\frac{\sin \theta_i}{2}
    +(1-p)e^{i\phi_j}\frac{\sin \theta_j}{2}
     &
    p\sin^2\frac{\theta_i}{2}
    +(1-p)\sin^2\frac{\theta_j}{2}
  \end{bmatrix}
\end{align*}

Looks complicated, so let's discuss the possible constraints.
From this, three constraint can be made.
Recall spherical coordinate
\begin{equation*}
  \mathbf r(x,y,z)
  =
  \mathbf r(\sin\theta\cos\phi,
  \sin\theta\sin\phi,\;
  \cos\theta)
\end{equation*}
First constraint is $\theta_2=-\theta_1=\theta$ and $ \phi_2=-\phi_1=\phi$
\begin{equation*}
  C_2=\begin{cases}
    x'
    =r\sin(-\theta)\cos(-\phi)
    =-r\sin\theta\cos\phi=-x, \\
    y'
    =r\sin(-\theta)\sin(-\phi)
    =r\sin\theta\sin\phi=y,   \\
    z'
    =r\cos(-\theta)=z
  \end{cases}
\end{equation*}
Next is $\theta_2=\theta_1$ and $ \phi_2=-\phi_1$
\begin{equation*}
  C_3=\begin{cases}
    x'
    =r\sin(\theta)\cos(-\phi)
    =r\sin\theta\cos\phi=x,   \\
    y'
    =r\sin(\theta)\sin(-\phi)
    =-r\sin\theta\sin\phi=-y, \\
    z'
    =r\cos(\theta)=z
  \end{cases}
\end{equation*}
Finally $\theta_2=-\theta_1$ and $ \phi_2=\phi_1$
\begin{equation*}
  C_4=\begin{cases}
    x'
    =r\sin(-\theta)\cos(\phi)
    =-r\sin\theta\cos\phi=-x, \\
    y'
    =r\sin(-\theta)\sin(\phi)
    =-r\sin\theta\sin\phi=-y, \\
    z'
    =r\cos(-\theta)=z
  \end{cases}
\end{equation*}

\subsubsection{Derivation: specific initial state.}
From the derived ergotropy
\begin{align*}
  \mathcal{E} \left[ \rho^{(1)}_{C_3} \right]
   & = B \left( 2\sin^2 \frac{\theta }{2} - 1 + \sqrt{1 - 4p(1-p)\sin ^2\theta \sin^2 \phi}  \right)
\end{align*}
we want to obtain initial state in which the coherence and population have the same contribution.
With $\theta > \pi/2$, we have
\begin{align*}
  4 \sin^2 \frac{\theta }{2 }-2             & =   1 - 2 \sin ^2 \frac{\theta }{2} + \sqrt{1-4p(1-p)\sin ^2\theta \sin^2 \phi} \\
  3\left( 2\sin^2\frac{\theta}{2}-1 \right) & = \sqrt{1-4p(1-p)\sin^2\theta\sin^2\phi}                                        \\
  - 3\cos\theta                             & =
  \sqrt{1-4p(1-p)\sin^2\theta\sin^2\phi}
\end{align*}
Since \(\theta>\pi/2\), the right-hand side is positive as required.
Writing it as $3 \cos \theta$ then squaring it
\begin{align*}
  9\cos^2\theta & =
  1-4p(1-p)\sin^2\theta\sin^2\phi                                                                           \\
  p^2-p         & =
  - \frac{1-9\cos^2\theta}{4\sin^2\theta\sin^2\phi}                                                         \\
  p             & = \frac{1 }{2 }\pm \frac{1 }{2} \sqrt{1 - \frac{1-9\cos^2\theta}{\sin^2\theta\sin^2\phi}}
\end{align*}
Since qubit ergotropy is symmetric around $p=1/2$, we can select the minus branch
\begin{equation*}
  p  = \frac{1 }{2 } - \frac{1 }{2} \sqrt{1 - \frac{1-9\cos^2\theta}{\sin^2\theta\sin^2\phi}}
\end{equation*}

Now the range $p \in [0,0.5]$, this  means that
\begin{align*}
  0\le
  \frac{1-9\cos^2\theta}
  {\sin^2\theta\sin^2\phi}
   & \le1
\end{align*}
and we have two inequality
\begin{equation*}
  1-9\cos^2\theta\ge0\qquad
  {1-9\cos^2\theta}\le
  {\sin^2\theta\sin^2\phi}
\end{equation*}
For the first inequality, we write
\begin{align*}
  |\cos\theta| & \le\frac13\implies
  \begin{cases}
    \cos \theta\le\frac13\implies \theta \in [\theta_0,2\pi] \\
    \cos\theta\ge-\frac13 \implies\theta \in [0,\pi-\theta_0]
  \end{cases}
\end{align*}
with $\theta_0 = \arccos(1/3)$ and $\theta_0 \approx 71^\circ$ numerically.
Combining both requirement and also imposing the population inversion requirement, we have
\begin{equation*}
  \pi/2 < \theta \le \pi - \theta_0
\end{equation*}

For the second inequality, we write
\begin{equation*}
  |\sin\phi|\ge
  \sqrt{  \frac{1-9\cos^2\theta}{\sin^2\theta}}
  =\sqrt{  \frac{\sin^2\theta-8\cos^2\theta}{\sin ^2\theta}}
  = \sqrt{ 1-8 \cot ^2 \theta}
\end{equation*}
The same applies to the second inequality
\begin{align*}
  \begin{cases}
    \sin \phi\ge  \sqrt{ 1-8 \cot ^2 \theta} \implies \phi\in  [\phi_0,\pi-\phi_0] \\
    \sin \phi \le -\sqrt{ 1-8 \cot ^2 \theta} \implies \phi\in [\pi+\phi_0,2\pi-\phi_0]
  \end{cases}
\end{align*}
with $\phi_0 = \arcsin (\sqrt{ 1-8 \cot ^2 \theta} )$.
The domain is then obtained by combining both of the requirement
\begin{equation*}
  \phi \in[\phi_0,\pi-\phi_0] \cup[\pi+\phi_0,2\pi-\phi_0]
\end{equation*}
Picking the first branch, we introduce another parameter to make $\phi$ dependent on $\theta$ as the case of $p$
\begin{equation*}
  \phi = \phi_0 + \lambda(\pi - 2\phi_0 )
\end{equation*}
with $0 \le \lambda \le 1$.
In particular, at $\lambda=1/2$, $\phi=\pi/2$ become completely independent of $\theta$ and
\begin{align*}
  p & =
  \frac{1 }{2 }- \frac{1 }{2} \sqrt{1 - \frac{1-9\cos^2\theta}{\sin^2\theta\sin^2\phi}}
  = \frac{1 }{2 }- \frac{1 }{2} \sqrt{8\cot^2\theta}                            \\
    & = \frac{1 }{2 }-  \sqrt{2}|\cot\theta|=\frac{1 }{2 }-  \sqrt{2}\cot\theta
\end{align*}

\subsubsection{Derivation: Specific final state.}
We consider two case here: final ergotropy such that contain only coherence, and population in which the initial state have both contribution.
This enforce
\begin{equation*}
  \sin ^2 \frac{\theta }{2 }>\frac{1 }{2} \implies \theta>2 \arcsin (1/\sqrt{2}) = \frac{\pi }{2}
\end{equation*}
The first case, coherence only, is trivial.
Population exist at receiver if
\begin{equation*}
  \sin ^2 \frac{\theta }{2 }|f_N|^2>\frac{1 }{2} \implies \theta>2 \arcsin \frac{1 }{\sqrt{2|f_N|^2}}
\end{equation*}
So the condition for coherence only transport
\begin{equation*}
  \frac{\pi }{2 }<\theta<2 \arcsin \frac{1 }{\sqrt{2|f_N|^2}}
\end{equation*}
The population only, meanwhile, is more complicated.
In terms of transfer amplitude, we require
\begin{equation*}
  |f_N|^2>\frac{1}{2\sin^2(\theta/2)}
\end{equation*}
Setting the coherence contribution to zero and imposing this condition
\begin{align*}
  \sqrt{1-4 |f_N|^2\left[p(1-p)\sin ^2\theta \sin^2\phi+\sin ^4\frac{\theta}{2}(1-|f_N|^2)\right]}
   & =
  2|f_N|^2\sin^2\frac{\theta}{2}-1
\end{align*}
Square it and simplify
\begin{multline*}
  1
  -4|f_N|^2p(1-p)\sin^2\theta\sin^2\phi
  -4|f_N|^2\sin^4\frac{\theta}{2}
  +4|f_N|^4\sin^4\frac{\theta}{2}
  \\
  =
  1
  -4|f_N|^2\sin^2\frac{\theta}{2}
  +4|f_N|^4\sin^4\frac{\theta}{2}.
\end{multline*}
to get
\begin{align*}
  -4|f_N|^2p(1-p)\sin^2\theta\sin^2\phi
       & = -4|f_N|^2\sin^2\frac{\theta}{2}+4|f_N|^2\sin^4\frac{\theta}{2}                      \\
  p(1-p)\sin^2\theta\sin^2\phi
       & = \sin^2\frac{\theta}{2}\left( 1- \sin ^2 \frac{\theta }{2} \right)                   \\
  p(1-p)\sin^2\theta\sin^2\phi
       & =  \frac{1 }{4 }\sin^2\theta                                                          \\
  \sin^2\phi
       & = \frac{1}{4p(1-p)}                                                                   \\
  \phi & = \begin{cases}
             \arcsin\sqrt{\frac{1}{4p(1-p)}} \\ \pi - \arcsin\sqrt{\frac{1}{4p(1-p)}}
           \end{cases}
\end{align*}
Real solution only exist for $p=0.5$ and yield $\phi=\pi/2,3\pi/2$, of which the coherence ergotropy doesn't exist in the first place.
So this is pointless.

\subsubsection{Derivation: ergotropy for useful state.}
Next $C_3(\theta_1=\theta_2,\phi_1=-\phi_2)$
\begin{align*}
  \rho^{(1)}_{C_3}  & =
  \begin{bmatrix}
    \cos ^2 \theta/2                                                               & \frac{\sin \theta  }{2} \left[p(e^{-i\phi}-e^{i\phi}) + e^{i\phi}  \right] \\
    \frac{\sin \theta  }{2} \left[p(e^{i\phi} - e^{-i\phi} ) + e^{-i\phi}  \right] & \sin ^2 \frac{\theta  }{2}
  \end{bmatrix} \\
  \rho ^{(1)}_{C_3} & =
  \begin{bmatrix}
    \cos ^2 \theta/2                                                 & \frac{\sin \theta}{2} \left[-i2p \sin \phi + e^{i\phi}  \right] \\
    \frac{\sin \theta  }{2} \left[i2p \sin\phi + e^{-i\phi}  \right] & \sin ^2 \frac{\theta  }{2}
  \end{bmatrix}
\end{align*}
and the final state evolution
\begin{equation*}
  \rho ^{(N)}_{C_3}  =
  \begin{bmatrix}
    1- \sin^2 \theta/2 |f_N|^2                                             & \frac{\sin \theta}{2} \left[-i2p \sin \phi + e^{i\phi}  \right]f_N \\
    \frac{\sin \theta  }{2} \left[i2p \sin\phi + e^{-i\phi}  \right]f_N ^* & \sin ^2 \frac{\theta  }{2}|f_N|^2
  \end{bmatrix}
\end{equation*}
The determinants are
\begin{align*}
  \det\rho^{(1)}_{C_3}
                         & =
  \sin^2\frac{\theta}{2}\cos^2\frac{\theta}{2}
  -
  \frac{\sin^2\theta}{4}
  \left| \cos \phi+ i(1-2p) \sin \phi\right|^2 \\
                         & =
  \frac{\sin^2\theta}{4}
  -\frac{\sin^2\theta}{4}
  \left[
    \cos^2\phi
    +(1-2p)^2\sin^2 \phi
  \right]                                      \\
                         & =
  \frac{\sin^2\theta}{4}
  \left[
    1-1-
    (4p^2-4p)\sin^2\phi
  \right]                                      \\
  \det \rho ^{(1)}_{C_3} & =
  p(1-p)\sin ^2\theta \sin^2 \phi
\end{align*}
and
\begin{equation*}
  \det \rho ^{(N)}_{C_3} =     |f_N|^2 \left[  p(1-p)\sin ^2\theta \sin^2 \phi + \sin^4 \frac{\theta}{2} (1-|f_N|^2)\right]
\end{equation*}
The initial ergotropy is
\begin{align*}
  \mathcal{E} \left[ \rho^{(1)}_{C_3} \right]
   & = B \left( 2\sin^2 \frac{\theta }{2} - 1 + \sqrt{1 - 4p(1-p)\sin ^2\theta \sin^2 \phi}  \right)
\end{align*}
and the final state
\begin{multline*}
  \mathcal{E} \left[ \rho^{(N)}_{C_3} \right]     = B \bigg( 2\sin^2 \frac{\theta }{2} |f_N|^2 - 1
  \\  \qquad+  \sqrt{1 - 4|f_N|^2 \left[ p(1-p)\sin ^2\theta \sin^2 \phi + \sin ^4 \frac{\theta }{2 }(1-|f_N|^2) \right]}  \bigg)
\end{multline*}
The population ergotropy for both
\begin{equation*}
  \mathcal{E}_\text{pop} \left[ \rho^{(1)}_{C_3} \right]= 2B(2 \sin^2 \frac{\theta }{2 }-1)  \qquad
  \mathcal{E}_\text{pop} \left[ \rho^{(N)}_{C_3} \right]= 2B(2 |f_N|^2\sin ^2\frac{\theta}{2 }-1)
\end{equation*}
And the coherence ergotropy
\begin{align*}
  \mathcal{E}_\text{coh}\left[ \rho^{(1)}_{C_3}\right]
   & = B \left[  1 - 2 \sin ^2 \frac{\theta }{2} + \sqrt{1-4p(1-p)\sin ^2\theta \sin^2 \phi}  \right]
\end{align*}
and
\begin{multline*}
  \mathcal{E}_\text{coh}\left[ \rho^{(N)}_{C_3}\right]
  = B \bigg[  1 - 2 |f_N|^2\sin ^2 \frac{\theta }{2} \\
    + \sqrt{1-4 |f_N|^2 \left[ p(1-p)\sin ^2\theta \sin^2 \phi
        + \sin ^4 \frac{\theta }{2 }(1-|f_N|^2) \right] }  \bigg]
\end{multline*}

\begin{figure}
  \centering
  \begin{tikzpicture}[tdplot_main_coords,scale=3.5]
    \coordinate (O) at (0,0,0);
    \def\r{1}
    \def\th{50}
    \def\ph{70}
    \def\phm{-\ph}
    \pgfmathsetmacro{\rxy }{sin(\th)}

    \draw[thick,->] (O) -- (1.1,0,0) node[anchor=north east]{$\hat{x}$};
    \draw[thick,->] (O) -- (0,1.1,0) node[anchor=north west]{$\hat{y}$};
    \draw[thick,->] (O) -- (0,0,1.1) node[anchor=south]{$\hat{z}$};
    \shade[ball color=white, opacity=0.1] (0,0,0) circle (1);

    \tdplotsetcoord{A}{\r}{\th}{\ph}
    \tdplotsetcoord{Ath}{\r}{\th/2}{\ph}
    \tdplotsetcoord{A2}{\r}{\th}{\phm}
    \tdplotsetcoord{A2th}{\r}{\th/2}{\phm}
    \coordinate (P) at ({sin(\th)*cos(\ph)},0,{cos(\th)});

    \draw[dashed] (A) -- node[midway,midway,xshift=-15pt,yshift=-8pt, left] {$1-p$}  node[midway,xshift=20pt,yshift=-12pt, right] {$ p$} (A2) ;

    \tdplotdrawarc[->]{(O)}{0.2}{0}{\ph}{anchor=north}{$\phi$}
    \tdplotdrawarc[->]{(O)}{0.2}{0}{\phm}{anchor=north east}{$-\phi$}
    \tdplotsetthetaplanecoords{\ph}
    \tdplotdrawarc[->,tdplot_rotated_coords]{(O)}{0.2}{0}{\th}{anchor=south west
    }{$\theta$}
    \tdplotsetthetaplanecoords{\phm}
    \tdplotdrawarc[->,tdplot_rotated_coords]{(O)}{0.2}{0}{\th}{anchor=south east}{$\theta$}

    \draw[dashed]  (O) -- (Axy);
    \draw[dashed]  (A) -- (Axy);
    \draw[dashed]  (O) -- (A2xy);
    \draw[dashed]  (A2) -- (A2xy);

    \begin{scope}[tdplot_main_coords]
      \pgfmathsetmacro{\gamma}{acos(sin(\th)*sin(\th)*cos(\ph-\phm)+cos(\th)*cos(\th))}
      \draw[dashed,domain=0:1,samples=100,smooth]
      plot (
      {sin((1-\x)*\gamma)/sin(\gamma)*sin(\th)*cos(\ph)+sin(\x*\gamma)/sin(\gamma)*sin(\th)*cos(\phm)},
      {sin((1-\x)*\gamma)/sin(\gamma)*sin(\th)*sin(\ph)+sin(\x*\gamma)/sin(\gamma)*sin(\th)*sin(\phm)},
      {cos(\th)}
      );
      \draw[dashed,domain=0:1,samples=100,smooth]
      plot (
      {{sin(\th)*cos(\ph)}},
      {sin((1-\x)*\gamma)/sin(\gamma)*sin(\th)*sin(\ph)+sin(\x*\gamma)/sin(\gamma)*sin(\th)*sin(\phm)},
      {sin((1-\x)*\gamma)/sin(\gamma)*cos(\th)+sin(\x*\gamma)/sin(\gamma)*cos(\th)}
      );
    \end{scope}

    \node[circle,fill=white] at (A) {};
    \node[circle,fill=white] at (A2) {};
    \node[circle,fill=white] at (P) {};
    \draw[thick,->,>=Stealth] (O) -- (A) node [anchor=south west]{$\ket{\psi_1}$};
    \draw[thick,->,>=Stealth] (O) -- (A2) node [anchor=south east]{$\ket{\psi_2}$};
    \draw[thick,->,>=Stealth] (O) -- (P) node [anchor=south ]{$\rho$};

    \def\alp{20}
    \def\bet{70}
    \tdplotsetrotatedcoords{\alp}{0}{0}
    \begin{scope}[tdplot_rotated_coords]
      \draw[thick] (0,-\r) arc (-90:90:\r);
      \draw[dashed] (0,\r) arc (90:270:\r);
    \end{scope}

    \tdplotsetrotatedcoords{\alp}{\bet}{0}
    \begin{scope}[tdplot_rotated_coords]
      \draw[thick] (0,\r) arc (90:270:\r);
      % \draw[thick] (\r,0) arc (0:360:\r);
    \end{scope}
  \end{tikzpicture}
  \caption*{Figure: Bloch sphere representation of mixed states created from pure state $\ket{\psi_i}$ ensemble in $C_3$.}
  \label{fig:chain}
\end{figure}

Take the third constraint $C_3$, the state in question is
\begin{align*}
  \rho^{(N)}
   & =
  \frac12
  \begin{bmatrix}
    1+z_N    & x_N-iy_N \\
    x_N+iy_N & 1-z_N
  \end{bmatrix} \\
  \rho^{(N)}
   & = \frac{1 }{2}
  \begin{bmatrix}
    2- 2|f_N|^2\sin^2 \theta/2                                 & \sin \theta \left[-i2p \sin \phi + e^{i\phi}  \right]f_N \\
    \sin \theta \left[i2p \sin\phi + e^{-i\phi}  \right]f_N ^* & 2|f_N|^2 \sin ^2 \theta/2
  \end{bmatrix}
\end{align*}
We can determine its final state Bloch vector.
Note that $f_N = |f_N|e^{i\delta}$, so the off diagonal matrices without prefactor can be written as
\begin{align*}
  \rho_{01}
   & =
  |f_N|\sin\theta\left[-i2p\sin\phi+e^{i\phi}\right]e^{i\delta}    \\
   & =
  |f_N|\sin\theta e^{i\delta}\left[\cos\phi+i(1-2p)\sin\phi\right] \\
   & =
  |f_N|\sin\theta
  \big\{\left[\cos\phi\cos\delta-(1-2p)\sin\phi\sin\delta\right]   \\
   &
  +i\left[\cos\phi\sin\delta+(1-2p)\sin\phi\cos\delta
    \right]\big\}
\end{align*}
While the diagonal
\begin{equation*}
  \rho_{11} =
  1-2\sin ^2 \frac{\theta  }{2}|f_N|^2 =
  1-|f_N|^2(1-\cos \theta)
\end{equation*}
Hence, the Bloch vector is
\begin{equation*}
  \begin{bmatrix}
    x_N \\
    y_N \\
    z_N
  \end{bmatrix}
  =
  \begin{bmatrix}
    |f_N|\sin\theta
    \left[
      \cos\phi\cos\delta
      -(1-2p)\sin\phi\sin\delta
    \right] \\
    |f_N|\sin\theta
    \left[
      (2p-1)\sin\phi\cos\delta
      -\cos\phi\sin\delta
    \right] \\
    1-|f_N|^2(1-\cos\theta).
  \end{bmatrix}
\end{equation*}
Let's try expressing in terms of initial vector.
Setting $|f_N|=1$, we can also obtain
\begin{align*}
  \rho_{01} & =  \sin \theta [\cos \phi + i(1-2p)\sin \phi]\qquad
  \rho_{11} = \sin ^2\theta/2
\end{align*}
and
\begin{equation*}
  \mathbf r^{(1)} =
  \begin{bmatrix} \sin\theta\cos\phi\\ (2p-1)\sin\theta\sin\phi\\ \cos\theta \end{bmatrix}
\end{equation*}
Then factoring receiver vector can then be written as
\begin{align*}
  \mathbf r^{(N)} & =
  \begin{bmatrix} |f_N|&0&0\\ 0&|f_N|&0\\ 0&0&|f_N|^2 \end{bmatrix}
  \begin{bmatrix}
    \cos\delta  & \sin\delta & 0 \\
    -\sin\delta & \cos\delta & 0 \\
    0           & 0          & 1
  \end{bmatrix}
  \begin{bmatrix} \sin\theta\cos\phi\\ (2p-1)\sin\theta\sin\phi\\ \cos\theta \end{bmatrix} \\
                  & \qquad+ \begin{bmatrix} 0\\ 0\\ 1-|f_N|^2 \end{bmatrix}
\end{align*}
or
\begin{equation*}
  \mathbf r^{(N)}  =  D(|f_N|)\,R_z(-\delta) \mathbf r^{(1)} +\mathbf t(|f_N|)
\end{equation*}
This is the affine transformation interpretation.
We can also interpret it in terms of initial state contribution
\begin{equation*}
  \mathbf r^{(N)}  =  M(|f_N|) \mathbf r^{(1)}_0 +(1-|f_N|^2)\mathbf r^{(N)}_0
\end{equation*}
The complex transfer amplitude $f_N = |f_N|e^{i\delta}$ introduce rotation about $z$ axis by angle $-\delta$, represented by
\begin{equation*}
  R_z (-\delta) =
  \begin{bmatrix}
    \cos\delta  & \sin\delta & 0 \\
    -\sin\delta & \cos\delta & 0 \\
    0           & 0          & 1
  \end{bmatrix}
\end{equation*}
The diagonal matrix is the anisotropic contraction produced by incomplete excitation transfer
\begin{equation*}
  D(|f_N|) = \operatorname{diag}(|f_N|,|f_N|,|f_N|^2)
\end{equation*}
where it act differently on transverse $(x,y)$ and longitudinal $(z)$ component, in general its action is to contracts the vector toward origin, thus reducing the purity.
The vector $\textbf{t}$ is an affine displacement along the \(+z\) direction
\begin{equation*}
  \textbf{t} = \begin{bmatrix} 0, 0,1-|f_N|^2 \end{bmatrix}
\end{equation*}
which add displacement toward the ground-state, thus restoring purity.

It is generally called the affine transformation of the Bloch vector (affine means linear transformation and translation).

We observe that at receiver site, coherence only ergotropy exist while population only doesn't.
Few reasons:
\begin{enumerate}
  \item $D(|f_N|)$ acts more strongly on $z$ axis, which is population-related component.
        Somehow related to the fact that off-diagonal coherence scales as $e^{-\lambda/2}$ and population as $e^{-\lambda}$?
  \item $\mathbf{t}(|f_N|)$ pushes \(z\) toward \(+1\), or $\ket{0}$.
        Since population ergotropy requires \(z<0\), this shift further suppresses population ergotropy.
        Somehow related to amplitude damping in an open system?
\end{enumerate}
The combination of the anisotropic contraction and the affine shift toward \(+z\) suppresses population inversion faster than transverse coherence, allowing coherence-only ergotropy but preventing population-only ergotropy for an initially coherent state.

We consider the imperfect transfer effect on stored energy $E_\mathrm{st} = E_0z$.
Therefore,
\begin{align*}
  E_{\rm st}^{(N)}
   & =  E_0 \left[ |f_N|^2z^{(1)} + 1-|f_N|^2 \right]       \\
  E_{\rm st}^{(N)}-E_{\rm st}^{(1)}
   & = E_0 \left[ |f_N|^2z^{(1)} +1-|f_N|^2-z^{(1)} \right] \\
  E_{\rm st}^{(N)}-E_{\rm st}^{(1)}
   & = E_0(1-|f_N|^2)(1-z^{(1)})\\
    E_{\rm st}^{(N)}-E_{\rm st}^{(1)} &\geq 0
\end{align*}
because the affine translation $\mathbf{t}$ moves the Bloch vector toward \(+z\), which corresponds to the ground state, we observe that  imperfect transfer always reduces the stored energy, except for the trivial cases \( |f_N|=1\) or \(z^{(1)}=1\).
Equivalently, it can be stated that during the evolution, increasing $|f_N|$ results in increasing  stored energy at receiver.

Due to $|\mathbf{r}|$ dependence, which can either increase or decrease depending on the initial state, the passive energy $E_\mathrm{pas}^{(N)} = E_0|r^{(N)}|$ and purity $\mathcal{P}^{(N)} = [1+|r^{(N)}|^2]/2$ is not universally monotonic.
We therefore consider the radial change of the Bloch vectors
\begin{align*}
  r_N^2
   & =  q(x_1^2+y_1^2)+[qz_1+(1-q)]^2                       \\
   & =  q(r_1^2-z_1^2)+[1-q(1-z_1)]^2                       \\
  r_N^2 -r_1^2
   & = qr_1^2-qz_1^2 + 1-2q(1-z_1) +q^2(1-z_1)^2-r_1^2      \\
   & = 1 - (1-q)r_1^2-qz_1^2 -2q(1-z_1) +q^2(1-z_1)^2       \\
   & = 1 - (1-q)r_1^2-q(1+(1+z_1)^2) +q^2(1-z_1)^2          \\
   & =  (1-q)(1-r_1^2)-q(1+z_1)^2 +q^2(1-z_1)^2             \\
  r_N^2-r_1^2
   & = (1-q)(1-r_1^2) - q(1-q)(1-z_1)^2                     \\
  r_N^2-r_1^2
   & =  (1-|f_N|^2)\left[ 1-r_1^2- |f_N|^2(1-z_1)^2 \right]
\end{align*}
with $q=|f_N|^2$, or equivalently
\begin{equation*}
  r_N-r_1 =
  \frac{  (1-|f_N|^2)\left[ 1-r_1^2- |f_N|^2(1-z_1)^2 \right]  }{r_N+r_1}
\end{equation*}
Since $0\leq |f_N|^2 \leq 1$, the sign is determined by the square braket.
The term $1-r_1^2$ measure the mixedness of the sender, while $|f_N|^2(1-z_1)^2$ measure the weighted displacement of the initial state from $\ket{0}$.

We define 
\begin{equation*}
 q_c= \frac{1-r_1^2}{(1-z_1)^2}\geq 0
\end{equation*}
as transfer probability at which the receiver has exactly the same purity as the sender.
And 
\begin{align*}
\frac{d r_N^2}{dq} 
& =  r_1^2-z_1^2-2(1-z_1) + 2q(1-z_1)^2 =0\\
q_m & = 
\frac{ 2(1-z_1)-(r_1^2-z_1^2) }{ 2(1-z_1)^2 } = \frac{1+q_c }{2} \geq \frac{1 }{2}
\end{align*}
as the transfer probability where the receiver has the minimum Bloch-vector length $r_N(q_m) = r_\text{min}$.
The first quantity is obtained by setting $r_N=r_1$, while the second is by setting its derivative instead.
Also 
\begin{equation*}
\frac{d^2 r_N^2}{dq^2} = 2(1-z_1)^2\ge0
\end{equation*}
Therefore \(r_N^2(q)\) is always convex $\cup$.
For $0\leq q \leq1$, we have $q_c\leq q_m$, hence the typical evolution, assuming monotonic increasing $|f_N|$ during transport, is
\begin{equation*}
  \underbrace{r_N(0)=1}_{\text{(a)}} \longrightarrow 
  \underbrace{r_N(q_c)=r_1}_{\text{(b)}} \longrightarrow 
  \underbrace{r_N(q_m)=r_{\min}}_\text{(c)} \longrightarrow 
  \underbrace{r_N(1)=r_1}_\text{(d)}
\end{equation*}
Do note that $r_N=r_1$ now refer to the condition (b) $q_c=q_m$ and (d) PST.

For $|f_N|^2<q_c$ or (a) to (b), receiver remains purer that the sender $r_N>r_1$ because the translation to $\ket{0}$ induced by $\mathbf{t}$ dominate.
Increasing $|f_N|^2$ weakens  $\mathbf{t}$ translation while  increasing sender contribution through reduced $D$ contraction, where the overall effect is to reduce $r_N$ until $r_N$ reaches $r_1$ at (b) and its minimum at (c).
After $q_m$, the transferred sender contribution becomes dominant, and the receiver increasingly resemble the sender.

By setting $q_m=q_c$ we see that it is only possible at $q=1$.
Then (b), (c), (d) coincide 
\begin{equation*}
   r_N(0)=1 \longrightarrow r_N(q_c)=r_N(q_m)= r_N(1) = r_1
\end{equation*}
Here the purity of receiver is the same as the sender, which is also receiver minimum.
This occur at initial state
\begin{align*}
r_1^2 = 1-(1-z_1)^2 = 2z_1-z^2  
\end{align*}
In our original parametrization, this reads 
\begin{align*}
\sin^2\theta\cos^2\phi +(2p-1)^2\sin^2\theta\sin^2\phi +\cos^2\theta & =  2\cos\theta-\cos^2\theta \\
 \sin^2\theta \left[ \cos^2\phi+(2p-1)^2\sin^2\phi \right] & =  2\cos\theta(1-\cos\theta) 
\end{align*}
For simplicity, choose $p=1/2$ and the condition reads 
\begin{align*}
\sin^2\theta\cos^2\phi & =  2\cos\theta(1-\cos\theta)\\
(1-\cos\theta)(1+\cos\theta)\cos^2\phi & =  2\cos\theta(1-\cos\theta)\\
 \cos^2\phi & =  
 \frac{2\cos\theta}{1+\cos\theta} 
\end{align*}
Hence we should observe that the radius decrease monotonically at 
\begin{equation*}
  p=\frac{1 }{2 },
  \qquad 0<\theta<\frac{\pi }{2 }, 
  \qquad \phi = \pm \arccos \sqrt{ \frac{2\cos\theta}{1+\cos\theta} }
\end{equation*}

Now suppose that initial state is pure $r_1=1$.
We have in this case $q_c=1$, thus (b) coincide with (a) and the evolution reads
\begin{equation*}
  r_N(0)=r_N(q_c)=1\longrightarrow
  r_N(q_m)=r_{\min}\longrightarrow
  r_N(1)=r_1
\end{equation*}



\subsubsection{Derivation: ergotropy for second mixed state (pointless).}
We impose then the constraint of $C_2(\theta_1=-\theta_2,\phi_1=-\phi_2)$, to get
\begin{align*}
  \rho^{(1)}_{C_2}  & =
  \begin{bmatrix}
    \cos ^2 \theta/2                                                               & \frac{\sin \theta  }{2} \left[p(e^{-i\phi}+e^{i\phi}) -e^{i\phi}  \right] \\
    \frac{\sin \theta  }{2} \left[p(e^{i\phi} + e^{-i\phi} ) - e^{-i\phi}  \right] & \sin ^2 \frac{\theta  }{2}
  \end{bmatrix} \\
  \rho ^{(1)}_{C_2} & =
  \begin{bmatrix}
    \cos ^2 \theta/2                                                & \frac{\sin \theta  }{2} \left[2p \cos \phi - e^{i\phi}  \right] \\
    \frac{\sin \theta  }{2} \left[2p \cos\phi - e^{-i\phi}  \right] & \sin ^2 \frac{\theta  }{2}
  \end{bmatrix}
\end{align*}
and the final state evolution
\begin{equation*}
  \rho ^{(N)}_{C_2} = \begin{bmatrix}
    1 - \sin^2 \frac{\theta }{2} |f_N|^2                                 & \frac{\sin \theta }{2} \left( 2p \cos \phi-e^{i\phi} \right)f_N \\
    \frac{\sin \theta }{2} \left( 2p \cos \phi-e^{-i\phi} \right) f_N ^* & \sin^2 \frac{\theta }{2 } |f_N|^2
  \end{bmatrix}
\end{equation*}
The determinants are
\begin{align*}
  \det\rho^{(1)}_{C_2}
                         & =
  \sin^2\frac{\theta}{2}\cos^2\frac{\theta}{2}
  -
  \frac{\sin^2\theta}{4}
  \left|(2p-1) \cos \phi- i \sin \phi\right|^2 \\
                         & =
  \frac{\sin^2\theta}{4}
  -\frac{\sin^2\theta}{4}
  \left[
    (2p-1)^2\cos^2\phi
    +\sin^2 \phi
  \right]                                      \\
                         & =
  \frac{\sin^2\theta}{4}
  \left[
    1-
    (4p^2-4p)\cos^2\phi-1
  \right]                                      \\
  \det \rho ^{(1)}_{C_2} & =
  p(1-p)\sin ^2\theta \cos^2 \phi
\end{align*}
and
\begin{align*}
  \det \rho^{(N)}_{C_2} & =
  |f_N|^2 \left[ p(1-p)\sin ^2\theta \cos^2 \phi + \sin ^4 \frac{\theta }{2 }(1-|f_N|^2) \right]
\end{align*}
Therefore the ergotropy is
\begin{align*}
  \mathcal{E} \left[ \rho^{(1)}_{C_2} \right] & =  B \left( 2\rho_{11}^{(1)}- 1 + \sqrt{1-4 \det \rho^{(1)}} \right)                            \\
                                              & = B \left( 2\sin^2 \frac{\theta }{2} - 1 + \sqrt{1 - 4p(1-p)\sin ^2\theta \cos^2 \phi}  \right)
\end{align*}
and
\begin{multline*}
  \mathcal{E} \left[ \rho^{(N)}_{C_2} \right]     = B \bigg( 2\sin^2 \frac{\theta }{2} |f_N|^2 - 1
  \\  \qquad+  \sqrt{1 - 4|f_N|^2 \left[ p(1-p)\sin ^2\theta \cos^2 \phi + \sin ^4 \frac{\theta }{2 }(1-|f_N|^2) \right]}  \bigg)
\end{multline*}
The population ergotropy for both
\begin{equation*}
  \mathcal{E}_\text{pop} \left[ \rho^{(1)}_{C_2} \right]= 2B(2 \sin^2 \frac{\theta }{2 }-1)  \qquad
  \mathcal{E}_\text{pop} \left[ \rho^{(N)}_{C_2} \right]= 2B(2 |f_N|^2\sin^2 \frac{\theta}{2 }-1)
\end{equation*}
And the coherence ergotropy
\begin{align*}
  \mathcal{E}_\text{coh}\left[ \rho^{(1)}_{C_2}\right]
   & = B \left[  1 - 2 \sin ^2 \frac{\theta }{2} + \sqrt{1-4p(1-p)\sin ^2\theta \cos^2 \phi}  \right]
\end{align*}
and
\begin{multline*}
  \mathcal{E}_\text{coh}\left[ \rho^{(N)}_{C_2}\right]
  = B \bigg[  1 - 2 |f_N|^2\sin ^2 \frac{\theta }{2} \\
    + \sqrt{1-4 |f_N|^2 \left[ p(1-p)\sin ^2\theta \cos^2 \phi
        + \sin ^4 \frac{\theta }{2 }(1-|f_N|^2) \right] }  \bigg]
\end{multline*}

\subsubsection{Derivation: ergotropy for fourth mixed state (even more pointless).}
Finally, the $C_4(\theta_1=-\theta_2,\phi_1=\phi_2)$ constraint
\begin{align*}
  \rho^{(1)}_{C_4}  & =
  \begin{bmatrix}
    \cos ^2 \theta/2                                                             & \frac{\sin \theta  }{2} \left[p(e^{-i\phi}+e^{-i\phi}) - e^{-i\phi}  \right] \\
    \frac{\sin \theta  }{2} \left[p(e^{i\phi} + e^{i\phi} ) - e^{i\phi}  \right] & \sin ^2 \frac{\theta  }{2}
  \end{bmatrix} \\
  \rho ^{(1)}_{C_4} & =
  \begin{bmatrix}
    \cos ^2 \theta/2                        & \frac{\sin \theta}{2}(2p-1) e^{-i\phi} \\
    \frac{\sin \theta  }{2} (2p-1)e^{i\phi} & \sin ^2 \frac{\theta  }{2}
  \end{bmatrix}
\end{align*}
and the final state evolution
\begin{equation*}
  \rho ^{(N)}_{C_4} =
  \begin{bmatrix}
    1- \sin^2 \theta/2|f_N|^2 & \frac{\sin \theta}{2}(2p-1) e^{-i\phi}f_N \\ \frac{\sin \theta  }{2} (2p-1)e^{i\phi}f_N ^*&  \sin ^2 \frac{\theta  }{2}|f_N|^2
  \end{bmatrix}
\end{equation*}
Not used, so I won't derive it further.

\subsubsection{Derivation: average ergotropy.}
In the same method as the fidelity derivation, we can determine the ergotropy transfer performance by averaging through all possible ergotropy value.
By defining $a = |f_N|^2$ and $s = \sin ^2 (\theta/2)$, we average
\begin{equation*}
  \mathcal{E}_\text{coh} (t) = 2as-1+\sqrt{1+4a(a-1)s^2}
\end{equation*}
over the solid angle
\begin{equation*}
  \braket{\mathcal{E}_\text{coh}} = \int d\Omega\; \mathcal{E}_\text{coh} = \frac{1 }{2 }\int_{0}^{\pi} d\theta\;\sin\theta \left[  2as-1+\sqrt{1+4a(a-1)s^2} \right]
\end{equation*}
with the factor $1/2$ due to $\mathcal{E}_\text{coh}$ being independent of $\phi$.
Substituting $ds = \sin (\theta/2)\cos (\theta/2)\;d\theta = \sin \theta\;d\theta/2$
\begin{equation*}
  \braket{\mathcal{E}_\text{coh}} = \int_{0}^{1} ds \left[  2as-1+\sqrt{1-4a(1-a)s^2} \right]
\end{equation*}
First and second term yield $a-1$, also note the sign changes on the third term.
With $4a(1-a)s^2=k^2s^2=\sin ^2\phi$ and $k\;ds=\cos\theta\;d\phi$, the third term reads
\begin{align*}
  \frac{1}{k }\int_{0}^{\arcsin k }d\phi\;\cos \phi \sqrt{1-\sin ^2\theta}
   & =
  \frac{1}{k }\int_{0}^{\arcsin k }d\phi\;\cos^2 \phi                                   \\
   & = \frac{1}{k }\int_{0}^{\arcsin k }d\phi\;\frac{1+\cos 2\phi}{2}                   \\
   & = \frac{1 }{k }\left[ \frac{\phi}{2}+ \frac{\sin 2\phi}{4} \right] _0^{\arcsin k } \\
   & = \frac{1 }{k }\left[ \frac{\arcsin k}{2}+ \frac{\sin 2\arcsin k}{4} \right]
\end{align*}
We write
\begin{equation*}
  \sin 2\arcsin k = 2 \sin(\arcsin k) \cos (\arcsin k) = 2k \cos \arcsin k
\end{equation*}
Since $k= \sin \theta$, $\cos \theta = \pm\sqrt{1+k^2}$, and $0\leq k\leq 1$ or equivalently $0\leq \arcsin k\leq \pi/2$ we have
\begin{equation*}
  \sin 2\arcsin k = 2k \sqrt{1-k^2} = 2k \sqrt{1-4a+4a^2} = 2k \sqrt{(1-2a)^2} = 2k|1-2a|
\end{equation*}
Therefore
\begin{equation*}
  \braket{\mathcal{E}_\text{coh}} = |f_N|^2-1 + \frac{\arcsin 2|f_N|\sqrt{1-|f_N|^2 }}{4|f_N|\sqrt{1-|f_N|^2}} + \frac{1 }{2} |1-2|f_N|^2 |
\end{equation*}

\subsubsection{Derivation: Purity.}
Using determinant obtained earlier, we plug into the general formula
\begin{align*}
  \mathrm{Tr}\left[\left(\rho^{(N)}_{C_2}\right)^2\right]
  =
  1
  -2|f_N|^2
  \left[
    p(1-p)\sin^2\theta\cos^2\phi
    +
    \sin^4\frac{\theta}{2}
    \left(1-|f_N|^2\right)
    \right]
\end{align*}
As sanity check, $p=1$ reduce into pure state case derived earlier
\begin{equation*}
  \mathrm{Tr}\left[\left(\rho^{(N)}_{C_2}\right)^2\right] =
  1
  -2\sin^4\frac{\theta}{2}
  |f_N|^2(1-|f_N|^2)
\end{equation*}
As for the third constraints
\begin{align*}
  \mathrm{Tr}\left[\left(\rho^{(N)}_{C_3}\right)^2\right]
  =
  1
  -2|f_N|^2
  \left[
    p(1-p)\sin^2\theta\sin^2\phi
    +
    \sin^4\frac{\theta}{2}
    \left(1-|f_N|^2\right)
    \right]
\end{align*}
\end{document}